\section{Marco conceptual y normativo}
\label{sec:marco_conceptual}

\subsection{Marco conceptual de tecnologías de información}

Para el desarrollo de la plataforma, se definen los siguientes conceptos básicos de desarrollo web, relacionados con mis objetivos de diseñar y desarrollar el sistema:

\begin{itemize}
    \item \textbf{Sistemas Web:} Aplicaciones que funcionan en el navegador y conectan al usuario con un servidor para realizar trámites de forma centralizada \parencite{pressman2020}.
    \item \textbf{Arquitectura MVC:} Forma de organizar la aplicación en tres partes: Modelo (datos y reglas), Vista (lo que se ve en pantalla) y Controlador (quien dirige las acciones), para mantener el sistema ordenado y fácil de mantener \parencite{pressman2020}.
    \item \textbf{HTML5:} Lenguaje que da estructura a las páginas y a los campos de los formularios \parencite{mdn2025}.
    \item \textbf{CSS3 y Diseño Responsivo:} Conjunto de estilos que adapta la página al tamaño de cualquier pantalla, ya sea celular o computadora \parencite{marcotte2018,w3schools2025}.
    \item \textbf{JavaScript:} Lenguaje que hace que la página responda al instante y avise si falta un dato antes de enviar el formulario \parencite{mozilla2025}.
    \item \textbf{PHP (Hypertext Preprocessor):} Lenguaje que trabaja en el servidor y se encarga de revisar los datos y guardarlos en la base de datos \parencite{arimetrics2024}.
    \item \textbf{PostgreSQL:} Base de datos libre y segura que guarda la información y cuida que los cupos no se dupliquen \parencite{postgresql2025}.
    \item \textbf{Autoservicio de Credenciales:} Opción que permite al estudiante recuperar solo su contraseña usando su carnet y fecha de nacimiento, sin ir a ventanilla.
\end{itemize}

\subsection{Marco conceptual del ámbito educativo institucional}

En el contexto del modelo educativo boliviano y de la institución sede, se conceptualizan los siguientes términos rectores:

\begin{itemize}
    \item \textbf{Bachillerato Técnico Humanístico (BTH):} Modelo formativo de la educación secundaria comunitaria productiva en Bolivia que otorga a los graduados una doble titulación: Bachiller Humanístico y Técnico Medio en una especialidad vocacional.
    \item \textbf{Módulo Tecnológico Productivo (MTP):} Centro educativo nucleador dotado de infraestructura de talleres y equipamiento industrial, que brinda formación técnica especializada a estudiantes provenientes de distintas unidades educativas del distrito municipal.
    \item \textbf{Preinscripción en Línea:} Proceso preliminar efectuado al cierre de la gestión previa mediante el cual el postulante reserva una vacante disponible en la especialidad técnica de su preferencia.
    \item \textbf{Matriculación Documental Definitiva:} Acto administrativo formal realizado al inicio de la gestión escolar correspondiente, en el cual se recibe, coteja y registra en el sistema la documentación obligatoria: libretas de calificaciones de 1ro a 3ro de secundaria, cédula de tutores, carnet de identidad y certificado de nacimiento.
    \item \textbf{Unidad Educativa de Procedencia (Colegio de Origen):} Establecimiento de educación secundaria regular al que pertenece formalmente el estudiante para su formación humanística, y que determina la compatibilidad de horarios y la conformación de grupos en el módulo.
    \item \textbf{Curso y Paralelo:} Grado académico (ej., 4to de secundaria) y sección de grupo (ej., A o B) en el que se organiza a los estudiantes para sus sesiones de taller según la capacidad áulica.
\end{itemize}

\subsection{Marco normativo boliviano}

Mi proyecto se apoya en las siguientes normas vigentes del Estado Plurinacional de Bolivia, relacionadas con mi objetivo de brindar un proceso justo y ordenado:

\begin{enumerate}
    \item \textbf{Constitución Política del Estado (CPE):} En sus Artículos 77 y 78 establece que la educación es la función más importante del Estado y debe ser para todos, productiva, gratuita y de calidad \parencite{cpe2009}.
    \item \textbf{Ley N° 070 de la Educación <<Avelino Siñani -- Elizardo Pérez>>:} Impulsa el Modelo Educativo Sociocomunitario Productivo y crea el Bachillerato Técnico Humanístico para unir el colegio con la formación técnica de cada región \parencite{ley070_2010}.
    \item \textbf{Resolución Ministerial RM 0912/2023:} Es el reglamento de graduación del BTH que permite presentar proyectos de <<Innovación Tecnológica>> para resolver problemas reales de los colegios y su comunidad \parencite{ministerioeducacion2023}.
\end{enumerate}
